\section{Collisions}
Now we will consider what happens to the momentum and energy of a test particle of charge $q_T$ and mass $m_T$ that is injected with velocity $v_T$ into a plasma. This test particle will make a sequence of random collisions with plasma particles (which we will call "field" particles); these collisions will alter both the momentum and energy of test particles.

The solution of the Rutherford scattering problem in the center of mass frame (see assignment 1 in chapter 1 of Bellan) shows that the scattering angle $\theta$ is given by
\begin{equation}
  \tan{\left(\frac{\theta}{2}\right)} = \frac{q_Tq_F}{4\pi\epsilon_0b\mu v_0^2} \sim \frac{\text{Coulomb interaction energy}}{\text{kinetic energy}},
\end{equation}

where $\mu^{-1} = m_T^{-1} + m_F^{-1}$ is the reduced mass, $b$ is the impact parameter, and $v_0$ is the initial relative velocity.

We'll separate collisions into two categories:
\begin{enumerate}
  \item large angle collisions where $\pi/2 \leq \theta \leq \pi$ and 
  \item small angle (grazing) collisions where $\theta \ll \pi/2$.
\end{enumerate}

This is defined somewhat arbitrarily, but doesn't actually matter too much, which will be discussed later.


\subsection{Proof of Small Angle Scattering Frequency}
For all the problems we've been considering, we start with Poisson's equation. So, let's take a moment to just look at Poisson's equation on its own.
\begin{eg}
  Tokamak
\end{eg}
\begin{explanation}
  something
\end{explanation}
